The Double Asteroid Redirection Test (DART) was a NASA mission to evaluate a
method of planetary defense against near-Earth objects. The spacecraft hit
Dimorphos, a moon of the asteroid Didymos, to study how the impact affected
its orbit.

\begin{description}
\item[(a)] Calculate the expected orbital period change (in minutes), assuming
  that the collision was head-on, central, and perfectly inelastic.

  Assume that before the impact Dimorphos orbited Didymos on a circular orbit
  with a period of $P = 11.92$\,h. The masses of Dimorphos and Didymos are
  $m = 4.3 \times 10^9$\,kg and $M = 5.6 \times 10^{11}$\,kg, respectively.
  The mass and speed of the DART spacecraft relative to Dimorphos at a moment
  of impact were $m_{\mathrm{s}} = 580$\,kg and
  $v_{\mathrm{s}} = 6.1$\,km\,s$^{-1}$. Neglect the gravitational influence of
  other bodies. \hfill(20 points)
\item[(b)] In reality, the orbital period of Dimorphos was observed to be
  changed by $\Delta P_0 = -33$\,min. This is due to the momentum transfer
  associated with the recoil of the ejected debris: the spacecraft was
  absorbed by the asteroid, but the impact excavated some material from the
  asteroid and ejected it into space. Calculate the momentum of the ejected
  debris and express it as a fraction of the momentum of Dimorphos before the
  collision. You can assume that the mass of the ejected material is much
  smaller than the mass of Dimorphos. \hfill(15 points)
\item[(c)] Calculate the velocity change (in mm\,s$^{-1}$) of Dimorphos as a
  result of the impact, taking into account the effect of the ejected debris.
  \hfill(5 points)
\end{description}
